\begin{textsample}{POS Dim 9 – global\_south\_2025\_05 – Score 2.00 – global\_south\_2025\_05/125012444.txt}  \label{ex:f9_pos_215}
Why do India and Israel stand alone?
THE tides of global diplomacy are shifting rapidly and two nations once seen as pillars of regional dominance -- India and Israel -- are finding themselves increasingly isolated on the world stage.
After India ' s recent war with Pakistan, the global community united in its condemnation of \textbf{Indian} aggression.
The overwhelming silence from the world in support of India, except a few isolated voices, exposed New Delhi ' s overestimation of its diplomatic capital.
Now, as Israel intensifies its brutal military campaign in Gaza and the West Bank, it too is heading down a path of growing international alienation.
On Monday, an unprecedented joint statement from British Prime Minister Keir Starmer, French President Emmanuel Macron and Canadian Prime Minister Mark Carney sharply rebuked the Israeli government.
The leaders declared their opposition to Israel ' s expansion of military operations in Gaza, condemned the intolerable human suffering and denounced the deliberate obstruction of humanitarian aid.
They warned Prime Minister Netanyahu that if Israel @ @ @ @ @ @ @ @ @ @"further concrete actions,"including targeted sanctions, would follow.
The statement also reaffirmed the illegality of permanent displacement and settlement expansions, calling them direct violations of international humanitarian law.
The call for an immediate ceasefire and the revival of efforts for a two-state solution echoed growing consensus from across the globe.
Yet, in this chorus of international outrage, one nation remains conspicuously silent -- India.
Far from condemning the assault, India appears to be tacitly encouraging Israel ' s savage military campaign, drawing disturbing parallels between their treatment of Palestinians and India ' s own conduct toward Pakistan and Kashmir.
Israeli Prime Minister Benjamin Netanyahu responded with fury, accusing the Western leaders of offering a"prize"for terrorism by pressuring Israel to halt its"defensive war."Rejecting the joint statement, Netanyahu claimed that calling for a Palestinian state rewarded Hamas for its October 7 attack and insisted that Israel would instead adhere to President Donald Trump ' s plan -- a controversial proposal involving the expulsion of @ @ @ @ @ @ @ @ @ @ the international community as ethnic cleansing.
Yet even Israel ' s closest allies are beginning to distance themselves.
Reports surfaced that Trump ' s envoys engaged in direct talks with Hamas for the release of American hostages without informing Tel Aviv -- a stark deviation from established protocol.
Israel was reportedly blindsided, learning of the negotiations through its own intelligence services, not Washington.
Similarly, Trump ' s recent decision to launch nuclear negotiations with Iran -- without prior notice to Netanyahu -- reflects a changing strategic calculus in Washington, even under a pro-Israel US Administration.
This growing defiance of Israeli overreach by traditional allies marks a turning point.
Just as India was humiliated by the international community after its failed military aggression against Pakistan, Israel now faces global condemnation for its indiscriminate bombing, blockade and mass displacement of Palestinians.
From targeted airstrikes on hospitals and schools to the systematic denial of food and medicine, Israel ' s actions have shocked the conscience of the world.
The clout of Gulf nations -- @ @ @ @ @ @ @ @ @ @ and their diplomatic and financial influence is shaping the future of the Middle East.
Turkey is asserting its regional role, Syria is re-emerging from isolation and Iran has weathered decades of sanctions to remain a formidable power.
The Arab world is no longer a passive observer ; it is an active participant in reshaping the region ' s geopolitical reality.
Countries across the Muslim world -- from Malaysia to Indonesia -- are rising as economic and military powers, building alliances and advocating for justice in Palestine.
The momentum for a just resolution to the Israeli-Palestinian conflict is accelerating.
The June 18 conference in New York, co-chaired by France and \textbf{Saudi} Arabia, will further solidify global consensus toward a two-state solution.
Even European countries long considered Israel ' s allies are now proposing recognition of a Palestinian state.
France has already hinted at such a move, while Canada and the UK are signalling that continued aggression will come at a diplomatic and economic cost.
Israel, much like India after its debacle @ @ @ @ @ @ @ @ @ @ isolation.
The message is clear : the world will no longer tolerate unchecked aggression disguised as self-defence.
The era of impunity is drawing to a close.
Israel must now choose between perpetuating its cycle of violence or joining the international community in earnest efforts for peace.
Its military dominance may buy time, but it can not buy legitimacy.
The writing is on the wall -- tyranny has a shelf life.
And if the Netanyahu government does not heed the warnings of its allies and adversaries alike, it risks leading Israel into an era of irreversible isolation and decline.
The Palestinian cause, long suppressed under the weight of occupation, bloodshed and despair, is finding renewed strength -- not just through armed resistance but through diplomacy, moral clarity and international unity.
Pakistan has shown how a smaller nation can defeat aggression with strategy and resilience.
Iran has shown how endurance and defiance can outlast decades of coercion.
The Middle East is changing and the world is @ @ @ @ @ @ @ @ @ @ might ; to coexistence, not conquest.
Just as India was forced to face the consequences of its arrogance, Israel too must prepare to be held accountable.
Liberation of Palestine is no longer a distant dream -- it is becoming an urgent global demand.

% matched lemmas: indian, saudi
\end{textsample}
